\section{The procedure inventory}
\label{sec:inventory}

Table~\ref{tab:inventory} sets out the seven steps. The remainder of this
section describes each: what the step is, what in it is capable of releasing
material, what is known, and what the evidence supports by way of control.

\begin{table*}[t]
\centering
\caption{The procedure inventory. Release is the inherent aerosol release score
on a 1--5 scale, before any control is credited. RBI is the residual biorisk
index on a 1--25 scale under baseline practice, shown both at the step's
primary agent profile and at the low-dose inhalational profile that a
tuberculosis or \textit{Brucella} specimen presents. Conditionality records
whether the evidence splits between the step as intended and the step with a
named failure or technique deviation.}
\label{tab:inventory}
\small
\begin{tabular}{@{}llccl cc@{}}
\toprule
 & \textbf{Bench step} & \textbf{Release} & \textbf{Evidence} &
 \textbf{Conditional on} & \textbf{RBI, primary} & \textbf{RBI, inhalational} \\
\midrule
S01 & Specimen reception and accessioning & 3.4 & \textbf{E2} & failure    & 5.67 & 15.87 \\
S02 & Centrifugation of primary tubes     & 3.4 & \textbf{E1} & failure    & 4.67 & 13.07 \\
S03 & Cap removal and aliquoting          & 4.0 & \textbf{E1} & technique  & 6.33 & 17.73 \\
S04 & Automated chemistry and haematology & 3.0 & \textbf{E5} & failure    & 3.33 & \phantom{0}9.33 \\
S05 & Cartridge-based molecular testing   & 3.6 & \textbf{E1} & technique  & 10.27 & 10.27 \\
S06 & MALDI-TOF identification            & 4.0 & \textbf{E3} & technique  & 17.73 & 17.73 \\
S07 & Open-tube and open-plate culture    & 4.6 & \textbf{E3} & ---        & 19.60 & 19.60 \\
\bottomrule
\end{tabular}
\end{table*}

\subsection{S01 Specimen reception and accessioning}

\textit{Boundary.} Opening transport bags and coolers, inspecting primary
containers, triaging leaking or broken containers, decanting into racks,
barcoding and sorting. Courier transport and outbound referral packaging are
outside.

No energy is applied deliberately at this step. The energy that matters arrives
already in the specimen, as a container that leaked in transit or breaks on
handling. Containment is therefore scored as mostly open rather than closed,
because leakage in transit is a routine rather than exceptional finding, and
because the worker leans over the container to read a label, placing the
breathing zone directly above any release.

No study has measured aerosol release during clinical specimen reception.
The step is graded \textbf{E2} on an inference from the simulated accident
measurements of Liu and colleagues, who recorded source intensities of roughly
$10^{3}$ to $10^{4}$~CFU\,m$^{-3}$ for a dropped container \cite{liu2023}. The
inference requires a stated caveat: that work used 30~mL of concentrated
bacterial suspension in a research biosafety laboratory, not a clinical
specimen in a receiving area, so the magnitude should not be read across
directly. What transfers is the ordering, not the number.

The step is failure-conditional, and the practical implication is that
controls belong at the breakage mode. A designated receipt station with a
bunded, non-porous surface and a leak-containment tray; a written trigger for
moving any visibly leaking or damaged container into a Class~II cabinet
unopened; rejection criteria that permit staff to refuse a container rather
than open it; and a spill kit within reach with a documented check. Respiratory
protection is not indicated for routine reception.

\subsection{S02 Centrifugation of primary specimen tubes}

\textit{Boundary.} Loading, balancing, running, and opening the lid and rotor
for capped primary tubes of blood, urine or body fluid. Centrifugation of open
vessels or of cultures is outside.

Centrifugation applies the highest mechanical energy routinely delivered to a
clinical specimen, and for that reason it is among the most frequently cited
aerosol-generating procedures in laboratory guidance. The measured evidence
does not support the fear in its usual form.

Dubuis and Duchaine air-sampled and surface-swabbed during centrifugation of
blood spiked with a model bacterium, in a working hospital biochemistry
laboratory, and recovered no organisms \cite{dubuis2021}. Two conditions must
travel with that result, and they are the reason the step is graded \textbf{E1}
but flagged: the centrifugation used closed-cap tubes in a swing-bucket rotor,
and the air sampling method had a detection limit of 4~CFU\,m$^{-3}$. A
negative result obtained with containment in place is evidence that the
containment works. It is not evidence that the bare operation is benign, and it
cannot exclude release below the detection limit.

The failure mode is separately attested. Ravnholdt and colleagues include
centrifuge spills among the accident types studied, and centrifuge mishaps
appear in the record of incidents associated with laboratory-acquired infection
\cite{ravnholdt2025}. Liu and colleagues measured centrifugal tube rupture at
roughly 10 to 100~CFU\,m$^{-3}$ \cite{liu2023}, which is real but, notably,
one to three orders of magnitude \textit{below} their measurements for dropped
containers and accidental injection. The common assumption that a centrifuge
accident is the worst thing that can happen in a pre-analytical area is not
supported by the only study that ranked them.

Controls follow the conditionality. Sealed safety buckets or a sealed rotor,
opened inside a Class~II cabinet whenever an inhalational hazard is plausible;
a lid interlock; documented inspection of rotors, buckets, gaskets and
O-rings; tube fill limits and a prohibition on spinning damaged tubes; and a
written breakage procedure. This brief deliberately specifies no numeric wait
interval before opening the lid, because no interval in the located evidence is
derived from measurement of a clinical centrifuge, and a number stated here
would acquire an authority the evidence does not support.

\subsection{S03 Cap removal and aliquoting}

\textit{Boundary.} Manual or automated removal of caps from evacuated primary
tubes, pouring or pipetting into secondary vessels, and recapping. Opening
positive blood culture bottles belongs to S07.

This is the only step in the inventory whose entire purpose is to remove
containment, and it scores accordingly on containment and on breathing-zone
proximity: hands rise to chest or face height and the operator looks directly
into the opening. A pressure differential is released, and the common practice
of rolling or flicking a stopper free adds energy at the worst moment.

Dubuis and Duchaine also sampled during the opening of spiked blood tubes and
recovered nothing \cite{dubuis2021}. The step is graded \textbf{E1} on that
result, with the same two conditions attached: the tubes were opened behind a
plexiglass panel, and the detection limit was 4~CFU\,m$^{-3}$. The step is
flagged technique-conditional because the measurement was obtained under
controlled removal; a stopper pulled with a snap, or a tube still warm and
under pressure from the centrifuge, has not been separately measured and should
not be assumed equivalent.

Controls in order of strength: closed-tube sampling wherever the analyser
supports it, which removes the step rather than controlling it; enclosed
automated decapping in preference to manual; manual opening of any specimen
flagged for mycobacterial, fungal or systemic bacterial culture performed
inside a Class~II cabinet; and a taught technique with the tube held away from
the face and the stopper covered with absorbent gauze, twisted rather than
pulled. A splash shield protects against splash and earns no credit against
inhalation.

\subsection{S04 Automated chemistry and haematology analysers}

\textit{Boundary.} Loading racks, cap-pierce and probe sampling, on-board
mixing and washing, body-fluid modes, clearing probe jams and short samples,
emptying liquid waste, and daily maintenance with covers open.

This is the highest-throughput step in the building. In a core laboratory it
runs effectively continuously, and a far larger number of specimens pass
through it than through any other step in this inventory.

It is graded \textbf{E5}. No located study has measured aerosol release from a
routine clinical chemistry or haematology analyser. The review by Ravnholdt and
colleagues states the gap directly, identifying highly automated laboratory
equipment as a source still requiring characterisation and noting that research
attention has concentrated on manual procedures and on accidents
\cite{ravnholdt2025}.

Normal running and fault clearance are, in substance, two different steps
sharing one name. In normal running the pathway is closed, with interfaces at
the pierced septum and the waste line. Essentially everything that could put
material into the air happens when a cover is open and a person is reaching
inside to clear a jam or perform maintenance. The step is therefore
failure-conditional, and a written jam and short-sample procedure that treats
cover-open work as a distinct task with its own controls is the single most
useful administrative control available.

Because no measurement exists, this brief recommends that a laboratory record
in its own risk assessment that it is operating on precaution rather than
evidence for this step. That record is itself an output required by a
management-system standard, and it is more honest than citing guidance that
does not exist.

\subsection{S05 Cartridge-based molecular testing}

\textit{Boundary.} Adding sample reagent to sputum or another primary specimen,
the required incubation, transferring treated specimen into the cartridge
chamber, closing the lid, loading the module, and disposing of the used
cartridge.

This step is treated in full in Section~\ref{sec:tb}, because its evidence is
the strongest in the inventory and is specific to tuberculosis work. In summary:
graded \textbf{E1} on direct measurement in which neither sample preparation
nor automated cartridge processing produced any culturable bioaerosol, against
16 and 325~CFU\,m$^{-3}$ measured during preparation of smears for acid-fast
microscopy \cite{banada2010}. The closed cartridge is a strong engineering
control that the laboratory receives without having to install anything, and
the brief's recommendation is to protect it rather than add to it.

\subsection{S06 MALDI-TOF target preparation and identification}

\textit{Boundary.} Picking colonies from a plate or transferring from positive
blood culture broth, smearing or extracting onto the target, applying matrix,
loading the target, and cleaning and reusing target plates. The primary culture
work that produced the colony belongs to S07.

The material handled here is among the most concentrated in the laboratory:
pure growth or enriched broth, worked close in under good lighting with the
face over an open target. Mechanical energy is low, but containment is absent
by design and concentration is high.

The step is graded \textbf{E3}. Ackelsberg and colleagues investigated
\textit{Brucella} exposure risk events across ten clinical laboratories in New
York City and found this identification technology in use during several of the
events; their recommendations included handling slow-growing organisms in a
cabinet and withholding automated identification until biological threat agents
have been excluded \cite{ackelsberg2020}. No measurement of aerosol release
during target preparation has been located.

The hazard sits upstream of the instrument, and the trigger that is supposed to
divert a dangerous isolate is where the failure occurs. In three of the four
events in which this technology featured, the initial Gram stain showed
Gram-positive rods or cocci suggesting streptococci. A control predicated on
the operator noticing a suspicious Gram stain would therefore have failed to
fire: the operators did look, and were misled. The control this brief
recommends is accordingly not a stain-based prompt but a blanket rule on
slow-growing blood culture isolates irrespective of stain appearance, combined
with a clinical and travel history cue, which is what the investigators
themselves recommended.

\subsection{S07 Open-tube and open-plate culture work}

\textit{Boundary.} Subculture, plate streaking and spreading, colony picking,
preparing suspensions, opening positive blood culture bottles, loop and needle
work, flaming, slide preparation from culture, and inspecting plates.

This step carries the highest release score and the highest residual biorisk in
the inventory, and it is the only step flagged as conditional on nothing: the
documented infections occurred during the step performed as intended, on an
open bench. It is graded \textbf{E3}, and the incident evidence is the
strongest in the brief.

Sejvar and colleagues identified sixteen cases of probable laboratory-acquired
meningococcal disease, of which eight were fatal; in fifteen of the sixteen
laboratories the isolate had been manipulated outside a biological safety
cabinet \cite{sejvar2005}. They estimated an attack rate among microbiologists
of 13 cases per 100,000 against 0.2 per 100,000 among adults in the general
population, an occupational rate roughly sixty-five times background, and
concluded that exposure to droplets or aerosols during open-bench manipulation
of invasive isolates was the most likely risk factor. The earlier MMWR
investigation records that affected microbiologists had performed subplating
with a loop on the open bench and had aspirated from blood culture bottles
without a cabinet, eye protection or a mask \cite{mmwr2002}. In a later fatal
case the investigation identified open-bench manipulation, open-bench plate
spreading, and flaming a slide that had not fully air dried \cite{mmwr2014}.
Across the \textit{Brucella} events, open-bench manipulation gave rise to more
than 200 occupational exposures requiring assessment and follow-up
\cite{ackelsberg2020}. Guidance additionally records cases arising from
smelling bacteriological cultures \cite{bmbl6}.

No aerosol measurement of routine subculture in a clinical laboratory has been
located. The step with the best-documented human harm in the inventory has
never had an air sampler placed beside it.

Cabinet use is accordingly the single highest-value control in this inventory
and the one carried by its strongest evidence: a certified Class~II cabinet for
all manipulation of sterile-site isolates and of any organism not yet ruled
out, with field certification on a defined schedule and recertification after
relocation. Disposable loops and a closed electric incinerator remove the need
to flame wet material. A standing prohibition on smelling plates should be
written down rather than assumed, and slides should be fully air dried before
heat is applied. Immunisation should be offered and recorded where an effective
vaccine exists for organisms staff routinely manipulate.
